\documentclass[12pt,oneside]{book}
\usepackage[a4paper,margin=2.5cm,headheight=25pt]{geometry}
\usepackage[T1]{fontenc}
\usepackage[utf8]{inputenc}
\usepackage[french]{babel}
\usepackage{lmodern}
\usepackage{microtype}
\usepackage{amsmath,amssymb}
\usepackage{graphicx}
\usepackage{booktabs,longtable,tabularx,array,multirow}
\usepackage{float}
\usepackage{caption}
\usepackage{xcolor}
\usepackage{enumitem}
\usepackage{fancyhdr}
\usepackage{titlesec}
\usepackage{tcolorbox}
\usepackage{lastpage}
\usepackage{hyperref}
\usepackage{bookmark}
\usepackage{csquotes}
\usepackage[backend=biber,style=authoryear,sorting=nyt]{biblatex}
\addbibresource{references.bib}

\definecolor{navy}{HTML}{12304A}
\definecolor{gold}{HTML}{B8903A}
\definecolor{lightnavy}{HTML}{EAF1F5}
\definecolor{darkgray}{HTML}{3F4A52}
\hypersetup{colorlinks=true,linkcolor=navy,citecolor=navy,urlcolor=navy,pdftitle={Segmentation intelligente des clients d'un distributeur en gros},pdfauthor={Saad Sahraoui et Hassan El Hajjam}}
\setlength{\parindent}{0pt}
\setlength{\parskip}{0.55em}
\renewcommand{\arraystretch}{1.2}
\setlist[itemize]{leftmargin=*,itemsep=2pt}
\setlist[enumerate]{leftmargin=*,itemsep=2pt}
\titleformat{\chapter}[display]{\bfseries\color{navy}}{\fontsize{15}{18}\selectfont\chaptername\ \thechapter}{8pt}{\Huge}
\titlespacing*{\chapter}{0pt}{-10pt}{24pt}
\titleformat{\section}{\Large\bfseries\color{navy}}{\thesection}{0.7em}{}
\titleformat{\subsection}{\large\bfseries\color{darkgray}}{\thesubsection}{0.7em}{}
\pagestyle{fancy}
\fancyhf{}
\lhead{\small\textcolor{navy}{Segmentation client — ML2}}
\rhead{\small\textcolor{darkgray}{ENSAM Rabat}}
\cfoot{\small\thepage\ / \pageref{LastPage}}
\renewcommand{\headrulewidth}{0.4pt}
\renewcommand{\headrule}{\hbox to\headwidth{\color{gold}\leaders\hrule height \headrulewidth\hfill}}
\newtcolorbox{pedagogie}{colback=lightnavy,colframe=navy,boxrule=0.5pt,arc=2pt,left=8pt,right=8pt,top=6pt,bottom=6pt,title=\textbf{Point pédagogique}}
\newtcolorbox{precaution}{colback=white,colframe=gold,boxrule=0.6pt,arc=2pt,left=8pt,right=8pt,top=6pt,bottom=6pt,title=\textbf{Précaution d'interprétation}}
\newcommand{\sourcefigure}{\small\textit{Source : élaboration propre à partir du dataset UCI Wholesale Customers.}}
\newcommand{\projectfigure}[3]{\begin{figure}[H]\centering\includegraphics[width=0.94\textwidth,height=0.60\textheight,keepaspectratio]{figures/#1}\caption{#2}\label{#3}\par\vspace{-2pt}\sourcefigure\end{figure}}
\newcommand{\smallfigure}[3]{\begin{figure}[H]\centering\includegraphics[width=0.82\textwidth,height=0.48\textheight,keepaspectratio]{figures/#1}\caption{#2}\label{#3}\par\vspace{-2pt}\sourcefigure\end{figure}}


\begin{document}
% La page de garde ne doit pas entrer en conflit avec l'ancre de la page 1
% du corps du rapport dans les signets PDF.
\hypersetup{pageanchor=false}
\begin{titlepage}
\centering
\vspace*{0.35cm}
\noindent
\begin{tabular*}{\textwidth}{@{\extracolsep{\fill}}cc@{}}
  \includegraphics[height=2.55cm,keepaspectratio]{figures/logo_um5.png} &
  \includegraphics[height=2.55cm,keepaspectratio]{figures/logo_ensam_rabat.png}
\end{tabular*}
\vspace{0.35cm}
{\Large\textsc{Université Mohammed V}\par}
\vspace{0.25cm}
{\large\textsc{École Nationale Supérieure d'Arts et Métiers — ENSAM Rabat}\par}
\vspace{0.9cm}
\vspace{0.35cm}
{\large Filière : Ingénierie Numérique en Data Science, Intelligence Artificielle et Santé Digitale (INDIA)\par}
\vspace{0.25cm}
{\large Module : Machine Learning 2\par}
\vfill
{\color{navy}\rule{\textwidth}{1.2pt}}\\[0.45cm]
{\Huge\bfseries\color{navy} Segmentation intelligente des clients\\[0.15cm] d'un distributeur en gros\par}
\vspace{0.25cm}
{\Large\textcolor{darkgray}{Étude comparative des algorithmes K-Means, CAH et DBSCAN}\par}
{\color{gold}\rule{\textwidth}{1.2pt}}
\vfill
\begin{tabular}{ll}
\textbf{Réalisé par :} & Saad Sahraoui et Hassan El Hajjam\\
\textbf{Encadré par :} & Pr. Safaâ Dafrallah\\
\textbf{Année universitaire :} & 2026/2027
\end{tabular}
\vfill
{\large\today\par}
\end{titlepage}

\frontmatter
\hypersetup{pageanchor=true}
\chapter*{Résumé}
Ce rapport présente une étude complète de segmentation non supervisée appliquée au dataset \textit{Wholesale Customers} de l'UCI Machine Learning Repository. L'objectif est d'identifier des groupes de clients d'un distributeur en gros à partir de six catégories de dépenses annuelles : \textit{Fresh}, \textit{Milk}, \textit{Grocery}, \textit{Frozen}, \textit{Detergents\_Paper} et \textit{Delicassen}. Les variables \textit{Channel} et \textit{Region}, bien que codées numériquement, sont exclues des distances et réservées à l'interprétation.

Après vérification de 440 observations sans valeur manquante, sans doublon exact et sans dépense négative, trois stratégies de préparation ont été étudiées. La stratégie principale, \texttt{log1p + StandardScaler}, réduit l'effet des asymétries positives et rend les amplitudes plus comparables. K-Means, la CAH Ward et DBSCAN sont ensuite comparés avec la silhouette, les effectifs, la couverture et l'interprétabilité. K-Means avec $K=2$ obtient une silhouette de 0,2903 et couvre l'ensemble des clients. La CAH Ward avec $K=2$ obtient 0,2585. DBSCAN révèle des zones denses mais classe 52,27\% des clients comme bruit dans la configuration analysée ; sa silhouette hors bruit n'est donc pas directement comparable. K-Means est retenu comme meilleur compromis parmi les configurations testées, tout en soulignant que la séparation reste modeste.

\textbf{Mots-clés :} apprentissage non supervisé, clustering, segmentation client, K-Means, CAH, DBSCAN, UCI.

\chapter*{Abstract}
This report presents a complete unsupervised customer-segmentation study based on the UCI \textit{Wholesale Customers} dataset. The analysis uses six annual spending variables and keeps the categorical variables \textit{Channel} and \textit{Region} for interpretation only. After data-quality checks and a comparison of three preprocessing strategies, the main representation is obtained with \texttt{log1p + StandardScaler}. K-Means, Ward agglomerative clustering and DBSCAN are evaluated using silhouette, cluster sizes, coverage and interpretability. K-Means with two clusters reaches a silhouette of 0.2903 on all 440 observations. Ward reaches 0.2585. DBSCAN identifies dense regions but labels 52.27\% of the observations as noise in the analysed configuration; its non-noise silhouette must therefore not be compared directly with full-dataset scores. K-Means is retained as the most balanced option among the tested configurations, although the separation remains modest.

\textbf{Keywords:} unsupervised learning, clustering, customer segmentation, K-Means, hierarchical clustering, DBSCAN.

\chapter*{Abréviations}
\begin{tabularx}{\textwidth}{>{\bfseries}lX}
ML & Machine Learning, apprentissage automatique\\
EDA & Exploratory Data Analysis, analyse exploratoire des données\\
CAH & Classification Ascendante Hiérarchique\\
PCA & Principal Component Analysis, analyse en composantes principales\\
IQR & Interquartile Range, écart interquartile\\
KDE & Kernel Density Estimate, estimation de densité par noyau\\
UCI & University of California, Irvine\\
\end{tabularx}

\tableofcontents
\listoffigures
\listoftables

\mainmatter
\chapter{Introduction générale}

\section{Contexte de la distribution en gros}
Les distributeurs en gros travaillent avec des clients dont les habitudes d'achat peuvent différer fortement. Un commerce de proximité, un hôtel ou un restaurant ne consomme pas nécessairement les mêmes catégories de produits, ni dans les mêmes proportions. Dans ce contexte, une lecture globale des ventes peut masquer des sous-populations importantes : certains clients achètent principalement des produits frais et surgelés, tandis que d'autres se concentrent sur l'épicerie, les produits laitiers ou les détergents.

La segmentation client consiste à regrouper des clients qui présentent des caractéristiques proches. Elle peut soutenir la description du portefeuille, la préparation d'offres différenciées et l'organisation d'analyses commerciales. Elle ne remplace toutefois pas la connaissance métier : un groupe mathématique n'est pas automatiquement un segment commercial validé.

\section{Limites d'une segmentation manuelle}
Une segmentation construite uniquement à partir de règles expertes peut être utile, mais elle dépend fortement des seuils choisis et devient difficile à maintenir lorsque les variables sont nombreuses. Les dépenses du dataset étudié ont des amplitudes très différentes et des distributions fortement asymétriques. Un seuil fixe peut alors favoriser les plus gros acheteurs ou confondre une différence de volume avec une différence de comportement.

L'apprentissage non supervisé offre une démarche complémentaire. Il permet de rechercher une structure dans les observations sans fournir à l'avance une classe correcte. Les groupes obtenus sont des hypothèses quantitatives à examiner, et non des vérités commerciales définitives.

\section{Problématique et question de recherche}
La question étudiée est la suivante :

\begin{pedagogie}
Comment identifier et interpréter des groupes de clients d'un distributeur en gros à partir de leurs dépenses annuelles, puis comparer K-Means, la CAH Ward et DBSCAN en tenant compte de la cohésion, de la séparation, de la couverture et de l'interprétabilité ?
\end{pedagogie}

Cette formulation impose deux précautions. Premièrement, les variables utilisées pour les distances doivent représenter les comportements de dépenses. Deuxièmement, les scores doivent être interprétés avec le contexte de chaque algorithme : DBSCAN peut laisser des observations en bruit, alors que K-Means et la CAH affectent tous les clients à un groupe.

\section{Objectifs du projet}
Les objectifs sont :
\begin{itemize}
  \item vérifier la structure et la qualité du dataset UCI Wholesale Customers ;
  \item décrire les distributions, les asymétries, les valeurs extrêmes et les corrélations ;
  \item comparer trois stratégies de préparation des six dépenses ;
  \item appliquer K-Means, CAH Ward et DBSCAN avec des paramètres explicitement documentés ;
  \item évaluer les partitions avec des mesures internes et les tailles des groupes ;
  \item revenir aux dépenses originales pour interpréter les profils ;
  \item formuler une conclusion prudente, sans présenter un cluster comme une causalité commerciale.
\end{itemize}

\section{Démarche générale}
La démarche suit une progression reproductible. Le fichier original est conservé dans \texttt{data/raw/}. Le contrôle initial vérifie les dimensions, les types, les valeurs manquantes et les doublons. L'EDA étudie les statistiques et les graphiques. Le prétraitement produit trois matrices distinctes. K-Means est ensuite étudié avec plusieurs valeurs de $K$, puis la CAH et DBSCAN sont exécutés sur la représentation principale \texttt{log1p + StandardScaler}.

Les résultats produits par les notebooks sont enregistrés dans \texttt{results/figures/} et \texttt{results/tables/}. Cette organisation sépare les données sources, les données préparées, le code d'analyse et les résultats consultables.

\section{Organisation du rapport}
Le chapitre 2 introduit les notions de clustering et de distance. Le chapitre 3 présente l'architecture technique et le protocole de reproductibilité. Le chapitre 4 décrit le dataset. Les chapitres 5 et 6 traitent respectivement l'EDA et le prétraitement. Les chapitres 7, 8 et 9 expliquent les trois algorithmes et leurs résultats. Le chapitre 10 réalise la comparaison, le chapitre 11 interprète les profils, et le chapitre 12 conclut l'étude. Les annexes regroupent les tableaux complémentaires et les repères pour la soutenance.

\chapter{Fondements théoriques}

\section{Machine Learning et apprentissage non supervisé}
Le Machine Learning désigne un ensemble de méthodes qui apprennent des régularités à partir de données. Dans un apprentissage supervisé, chaque observation est associée à une cible connue : par exemple une classe ou une valeur à prédire. Dans un apprentissage non supervisé, aucune cible de ce type n'est fournie. L'algorithme cherche une organisation interne des observations.

Dans ce projet, le fichier ne contient pas de label de cluster. \texttt{Channel} et \texttt{Region} décrivent des catégories observées, mais elles ne constituent pas une vérité cible pour entraîner les distances. Les clusters sont donc découverts uniquement à partir des six dépenses.

\begin{pedagogie}
Une classe est une catégorie connue ou définie par une règle externe. Un cluster est un groupe découvert par l'algorithme selon une notion de proximité ou de densité. Deux clusters portant le même numéro dans deux méthodes ne représentent pas nécessairement les mêmes clients.
\end{pedagogie}

\section{Observations, variables et features}
Une observation correspond ici à un client. Une variable est une colonne mesurée pour chaque client. Le terme \textit{feature} désigne une variable fournie à l'algorithme pour construire les distances. Les six features retenues sont \texttt{Fresh}, \texttt{Milk}, \texttt{Grocery}, \texttt{Frozen}, \texttt{Detergents\_Paper} et \texttt{Delicassen}.

Un label de cluster est une sortie de l'algorithme. Il ne doit pas être confondu avec une feature d'entrée. \texttt{customer\_index} est un identifiant de ligne ; l'ajouter aux distances introduirait une proximité artificielle liée à l'ordre du fichier.

\section{Distances, cohésion et séparation}
Pour deux vecteurs $x=(x_1,\ldots,x_p)$ et $y=(y_1,\ldots,y_p)$, la distance euclidienne est :
\begin{equation}
d(x,y)=\sqrt{\sum_{j=1}^{p}(x_j-y_j)^2}.
\end{equation}
La cohésion mesure si les observations d'un même groupe sont proches. La séparation mesure si les groupes sont éloignés entre eux. Ces notions dépendent de l'espace utilisé : une feature avec une amplitude beaucoup plus grande peut dominer la distance si elle n'est pas préparée.

La fonction d'inertie de K-Means est un exemple de critère de cohésion :
\begin{equation}
J=\sum_{k=1}^{K}\sum_{x_i\in C_k}\lVert x_i-\mu_k\rVert^2,
\end{equation}
où $C_k$ est le cluster $k$ et $\mu_k$ son centroïde. La séparation n'est pas une propriété absolue : un score doit être lu avec le nombre de groupes, les effectifs et la représentation choisie.

\section{Familles d'algorithmes}
Les méthodes de partition comme K-Means recherchent directement plusieurs groupes. Les méthodes hiérarchiques construisent une succession de fusions ou de divisions, ce qui permet de visualiser une structure multi-niveaux. Les méthodes de densité recherchent des zones où les observations sont suffisamment nombreuses et peuvent reconnaître du bruit.

K-Means et la CAH Ward utilisent une logique de groupes compacts dans l'espace préparé. DBSCAN utilise plutôt la densité locale. Cette différence explique pourquoi leurs résultats ne doivent pas être comparés uniquement avec un nombre décimal : leurs objectifs et leurs ensembles d'observations évalués peuvent différer.

\section{Score de silhouette}
Pour une observation $i$, on note $a(i)$ sa distance moyenne aux observations de son propre cluster et $b(i)$ la plus petite distance moyenne à un autre cluster. La silhouette individuelle est :
\begin{equation}
s(i)=\frac{b(i)-a(i)}{\max(a(i),b(i))}.
\end{equation}
Le score moyen est la moyenne de $s(i)$. Une valeur proche de 1 suggère des groupes bien séparés, une valeur proche de 0 un chevauchement, et une valeur négative une affectation possiblement mal adaptée. Cette lecture est indicative : la silhouette dépend de la distance, du prétraitement et de la partition.

\section{Pourquoi préparer les données ?}
Les six dépenses sont positives, asymétriques et d'amplitudes différentes. La transformation logarithmique \texttt{log1p} est définie par $\log(1+x)$ ; elle comprime les grandes valeurs tout en restant définie pour zéro. Une standardisation transforme ensuite chaque feature selon sa moyenne et son écart-type. Ces opérations ne changent pas le nombre de clients, mais changent la géométrie utilisée par les algorithmes.

\chapter{Architecture technique du projet}

\section{Environnement logiciel}
Le projet est développé sous Windows avec Python 3.12, VS Code et Jupyter. Les principales bibliothèques sont Pandas et NumPy pour la manipulation numérique, Matplotlib et Seaborn pour les figures, Scikit-learn pour le prétraitement et les algorithmes, et SciPy pour le dendrogramme hiérarchique. Git et GitHub assurent le versionnement du projet.

Codex a servi d'assistance au développement et à la documentation. Il ne remplace pas la vérification des sorties : les résultats présentés dans ce rapport proviennent des fichiers exécutés et des tableaux exportés dans le dépôt.

\section{Organisation des fichiers}
La structure complète est documentée dans le README du projet. Les répertoires ont les responsabilités suivantes :
\begin{description}
  \item[\texttt{data/raw/}] conserve le CSV original sans modification ;
  \item[\texttt{data/processed/}] contient les trois matrices préparées et le fichier des transformateurs ;
  \item[\texttt{notebooks/}] contient cinq notebooks exécutés dans l'ordre ;
  \item[\texttt{src/}] regroupe le contrôle du dataset et les fonctions de prétraitement ;
  \item[\texttt{models/}] conserve le modèle K-Means final ;
  \item[\texttt{results/figures/}] contient les 18 graphiques exportés ;
  \item[\texttt{results/tables/}] contient les métriques, les labels, les centroïdes et les profils ;
  \item[\texttt{reports/}] accueille le présent projet LaTeX.
\end{description}

\section{Pipeline reproductible}
Le pipeline conceptuel est :
\begin{center}
\fbox{Dataset} $\rightarrow$ \fbox{Analyse} $\rightarrow$ \fbox{EDA} $\rightarrow$ \fbox{Prétraitement} $\rightarrow$ \fbox{Clustering} $\rightarrow$ \fbox{Évaluation} $\rightarrow$ \fbox{Interprétation}
\end{center}

Le notebook 01 contrôle la structure. Le notebook 02 produit l'EDA. Le notebook 03 compare les stratégies de prétraitement et exporte les matrices. Le notebook 04 entraîne K-Means avec \texttt{random\_state=42}. Le notebook 05 compare la CAH Ward et DBSCAN aux résultats K-Means existants.

\section{Reproductibilité et traçabilité}
La reproductibilité repose sur plusieurs choix : les chemins sont construits à partir de la racine du projet, les features sont conservées dans un ordre explicite, les graines aléatoires sont fixées pour K-Means, et les résultats sont exportés avec des noms stables. Le CSV original n'est jamais écrasé.

Les données préparées utilisent un index \texttt{customer\_index}, qui est uniquement un identifiant. Les labels sont réalignés avec le fichier original par ordre de ligne avant de calculer les profils de dépenses. Cette étape évite d'attribuer un profil au mauvais client.

\chapter{Présentation du dataset}

\section{Origine et périmètre}
Le dataset \textit{Wholesale Customers} est disponible dans l'UCI Machine Learning Repository sous l'identifiant 292 \cite{uciWholesale}. Il comporte 440 observations et huit colonnes. Chaque ligne correspond à un client et les variables de dépenses représentent des montants annuels dont l'unité monétaire n'est pas précisée par la source.

Le fichier contient six dépenses quantitatives et deux variables codées : \texttt{Channel} et \texttt{Region}. Ces deux dernières sont utiles pour décrire les groupes après clustering, mais leur codage entier ne signifie pas qu'elles sont des mesures continues. Les distances initiales sont donc calculées uniquement sur les six dépenses.

\section{Description des variables}
\begin{table}[H]
\centering\small
\caption{Description des huit variables du dataset.}
\label{tab:variables}
\begin{tabularx}{\textwidth}{>{\ttfamily}l X l X}
\toprule
Variable & Signification & Nature & Rôle \\
\midrule
Channel & Canal commercial codé & Catégorielle & Interprétation \\
Region & Région codée & Catégorielle & Interprétation \\
Fresh & Produits frais & Quantitative & Feature de clustering \\
Milk & Produits laitiers & Quantitative & Feature de clustering \\
Grocery & Épicerie & Quantitative & Feature de clustering \\
Frozen & Produits surgelés & Quantitative & Feature de clustering \\
Detergents\_Paper & Détergents et papier & Quantitative & Feature de clustering \\
Delicassen & Produits delicatessen & Quantitative & Feature de clustering \\
\bottomrule
\end{tabularx}
\end{table}

\section{Premières observations et choix des features}
Les dépenses sont toutes strictement positives dans le fichier : aucune valeur nulle n'a été observée pour les six catégories. Les valeurs ne sont pas homogènes : \texttt{Fresh} atteint 112151 alors que \texttt{Delicassen} atteint 47943, et les écarts-types sont eux aussi très différents. Un algorithme fondé sur les distances pourrait donc donner trop de poids aux variables ayant les plus grandes amplitudes.

\texttt{Channel} possède deux modalités, 1 et 2, et \texttt{Region} trois modalités, 1, 2 et 3. Leur présence dans un tableau Pandas en \texttt{int64} ne les transforme pas en variables quantitatives adaptées à une distance euclidienne. Les inclure directement créerait une hiérarchie numérique artificielle entre les codes.

\begin{precaution}
Les labels de clusters n'existent pas dans le dataset. Les regroupements étudiés sont des structures découvertes par les algorithmes ; ils devront être interprétés comme des profils exploratoires et non comme des catégories officielles de clients.
\end{precaution}

\section{Qualité structurelle initiale}
Le contrôle effectué dans \texttt{src/check\_dataset.py} confirme 440 lignes, huit colonnes, zéro valeur manquante et zéro doublon exact. Les six dépenses ne contiennent aucune valeur négative. Les résultats détaillés de l'EDA sont présentés au chapitre suivant.

\chapter{Analyse exploratoire des données}

\section{Objectif de l'EDA}
L'analyse exploratoire cherche à comprendre les données avant tout calcul de distance. Elle vérifie la qualité du fichier, décrit la dispersion, mesure l'asymétrie, repère les valeurs extrêmes potentielles et examine les associations entre dépenses. Les graphiques ont été produits dans le notebook 02 avec Pandas, NumPy, Matplotlib et Seaborn.

\section{Statistiques descriptives}
\begin{table}[H]
\centering\scriptsize
\caption{Statistiques descriptives des six dépenses originales.}
\label{tab:descriptives}
\begin{tabular}{lrrrrrrr}
\toprule
Variable & Moyenne & Médiane & Écart-type & Min. & Q1 & Q3 & Max. \\
\midrule
Fresh & 12000,30 & 8504,00 & 12647,33 & 3 & 3127,75 & 16933,75 & 112151 \\
Milk & 5796,27 & 3627,00 & 7380,38 & 55 & 1533,00 & 7190,25 & 73498 \\
Grocery & 7951,28 & 4755,50 & 9503,16 & 3 & 2153,00 & 10655,75 & 92780 \\
Frozen & 3071,93 & 1526,00 & 4854,67 & 25 & 742,25 & 3554,25 & 60869 \\
Detergents\_Paper & 2881,49 & 816,50 & 4767,85 & 3 & 256,75 & 3922,00 & 40827 \\
Delicassen & 1524,87 & 965,50 & 2820,11 & 3 & 408,25 & 1820,25 & 47943 \\
\bottomrule
\end{tabular}
\end{table}

Pour chaque variable, la moyenne dépasse la médiane. Cette différence est cohérente avec une queue à droite : quelques clients ont des dépenses très supérieures au niveau central. Le coefficient de variation, défini par $CV=\sigma/\bar{x}$, dépasse 1 pour Milk, Grocery, Frozen, Detergents\_Paper et Delicassen ; il est proche de 1,05 pour Fresh. Cette dispersion justifie de ne pas utiliser les montants bruts directement dans les distances.

\section{Qualité des données}
Les contrôles donnent zéro valeur manquante, zéro doublon exact et zéro dépense négative. Aucune dépense n'est égale à zéro dans les six colonnes, ce qui signifie que l'étude de \texttt{log1p} n'est pas motivée par la présence de zéros, mais par la compression des grandes valeurs et la stabilité numérique de la formule.

\section{Distributions et asymétrie}
La figure \ref{fig:hist} rassemble les six histogrammes et les KDE. Elle montre des distributions concentrées sur des valeurs relativement faibles et prolongées par des queues droites. Les coefficients de skewness sont : Fresh 2,561, Milk 4,054, Grocery 3,587, Frozen 5,908, Detergents\_Paper 3,632 et Delicassen 11,152. Delicassen et Frozen sont donc les variables les plus asymétriques.

\projectfigure{01_histogrammes_kde_depenses.png}{Histogrammes et estimations KDE des dépenses.}{fig:hist}

La KDE est une représentation lissée de la densité et ne remplace pas les observations individuelles. Ici, l'intérêt principal est de visualiser les queues et la concentration, avant de choisir une transformation.

\section{Valeurs extrêmes et règle IQR}
Pour une variable donnée, l'écart interquartile est $IQR=Q_3-Q_1$. Les bornes conventionnelles sont :
\begin{equation}
[Q_1-1,5\,IQR,\;Q_3+1,5\,IQR].
\end{equation}
Les observations situées en dehors de ces bornes sont signalées, sans être automatiquement supprimées. Les nombres de signalements sont Fresh 20, Milk 28, Grocery 24, Frozen 43, Detergents\_Paper 30 et Delicassen 27, soit 172 signalements cumulés. Un même client peut être signalé dans plusieurs colonnes.

\projectfigure{02_boxplots_depenses_iqr.png}{Boxplots et valeurs signalées par la règle IQR.}{fig:boxplots}

Ces signalements peuvent correspondre à de gros acheteurs réels. Les supprimer sans vérification métier pourrait éliminer précisément les profils les plus intéressants. Le prétraitement utilise donc \texttt{log1p} et des mises à l'échelle, mais ne retire aucune ligne.

\section{Corrélations}
La corrélation de Pearson mesure une association linéaire standardisée. La plus forte est observée entre Grocery et Detergents\_Paper ($r=0,925$), suivie de Milk--Grocery ($r=0,728$) et Milk--Detergents\_Paper ($r=0,662$). La corrélation négative la plus marquée est Frozen--Detergents\_Paper ($r=-0,132$), donc faible en valeur absolue.

\projectfigure{03_heatmap_correlations_depenses.png}{Matrice de corrélation de Pearson des six dépenses.}{fig:corr}

La forte association Grocery--Detergents\_Paper signifie que, dans cet échantillon, les clients qui dépensent davantage dans une catégorie tendent aussi à dépenser davantage dans l'autre. Elle ne démontre pas que l'une cause l'autre. Elle indique toutefois que ces deux dimensions peuvent contribuer à distinguer un profil orienté épicerie et produits d'entretien.

\section{Relations bivariées et variables catégorielles}
Les nuages de points de la figure \ref{fig:scatter} utilisent \texttt{Channel} comme couleur descriptive. La couleur aide à observer des comportements différents entre canaux, mais \texttt{Channel} n'est pas une feature d'apprentissage.

\projectfigure{04_scatterplots_relations_depenses_channel.png}{Relations bivariées entre dépenses, colorées selon Channel.}{fig:scatter}

Les répartitions de Channel et Region sont présentées figure \ref{fig:categories}. Channel comprend 298 clients dans la modalité 1 et 142 dans la modalité 2. Region comprend 77 clients dans la modalité 1, 47 dans la modalité 2 et 316 dans la modalité 3. La table de contingence est :
\begin{table}[H]
\centering
\caption{Table de contingence Channel par Region.}
\label{tab:contingence}
\begin{tabular}{lrrr}
\toprule
 & Region 1 & Region 2 & Region 3 \\
\midrule
Channel 1 & 59 & 28 & 211 \\
Channel 2 & 18 & 19 & 105 \\
\bottomrule
\end{tabular}
\end{table}

\projectfigure{05_repartition_channel_region.png}{Répartition de Channel, de Region et table visuelle de contingence.}{fig:categories}

\section{Dépenses par Channel}
Les moyennes montrent une différence descriptive nette. Le canal 1 a une moyenne Fresh de 13475,56 et Frozen de 3748,25, tandis que le canal 2 a une moyenne Milk de 10716,50, Grocery de 16322,85 et Detergents\_Paper de 7269,51. Les médianes suivent la même direction générale : le canal 2 a une médiane Grocery de 12390 contre 2684 pour le canal 1. Ces écarts peuvent évoquer des portefeuilles commerciaux différents, mais ne prouvent pas leur cause.

\projectfigure{06_comparaison_depenses_par_channel.png}{Comparaison des dépenses par Channel.}{fig:channel}

\section{Synthèse EDA}
L'EDA conduit à quatre décisions : sélectionner les six dépenses, exclure Channel et Region des distances initiales, ne pas supprimer les signalements IQR et transformer les montants avant clustering. La skewness positive et les amplitudes hétérogènes rendent \texttt{log1p + StandardScaler} particulièrement pertinente comme représentation principale, tout en conservant StandardScaler seul et \texttt{log1p + RobustScaler} pour la sensibilité.

\chapter{Prétraitement des données}

\section{Pourquoi prétraiter ?}
Les méthodes fondées sur les distances utilisent directement la géométrie des features. Une variable dont les valeurs sont plus grandes peut dominer la distance euclidienne, même si elle n'est pas la plus informative pour distinguer les clients. De plus, les distributions très asymétriques rendent les moyennes et les centroïdes sensibles à quelques grandes valeurs.

Le prétraitement est donc une décision de modélisation. Il ne doit pas être confondu avec un nettoyage qui supprimerait les observations : les 440 clients sont conservés dans toutes les matrices.

\section{Sélection des variables}
Les six colonnes \texttt{Fresh}, \texttt{Milk}, \texttt{Grocery}, \texttt{Frozen}, \texttt{Detergents\_Paper} et \texttt{Delicassen} sont extraites dans un ordre fixe. \texttt{Channel} et \texttt{Region} sont conservées dans le fichier original pour les profils, mais ne sont pas passées aux algorithmes. L'index \texttt{customer\_index} sert à réaligner les lignes et n'est jamais une feature.

\section{Transformations utilisées}
La transformation logarithmique est :
\begin{equation}
x' = \log(1+x).
\end{equation}
Le terme 1 rend la transformation définie pour $x=0$. Dans ce dataset, les dépenses sont positives, mais \texttt{log1p} reste la fonction réutilisable documentée dans le module Python.

La standardisation utilise :
\begin{equation}
z=\frac{x-\mu}{\sigma},
\end{equation}
où $\mu$ et $\sigma$ sont calculés sur la feature d'entraînement. Le résultat est centré et réduit à la précision numérique près.

Le RobustScaler utilise la médiane et l'écart interquartile :
\begin{equation}
r=\frac{x-\operatorname{mediane}(x)}{IQR(x)}.
\end{equation}
Il est moins sensible aux valeurs extrêmes pour le centrage et l'échelle, mais il ne corrige pas à lui seul la skewness.

\section{Trois stratégies comparées}
\begin{table}[H]
\centering
\caption{Stratégies de prétraitement étudiées.}
\label{tab:strategies}
\begin{tabularx}{\textwidth}{>{\bfseries}l X X}
\toprule
Stratégie & Opérations & Lecture \\
\midrule
A & StandardScaler seul & Met les variables sur une échelle comparable, sans réduire directement l'asymétrie. \\
B & log1p puis StandardScaler & Compresse les grandes valeurs puis standardise ; candidate principale. \\
C & log1p puis RobustScaler & Compresse les grandes valeurs et utilise médiane/IQR ; alternative de sensibilité. \\
\bottomrule
\end{tabularx}
\end{table}

\projectfigure{07_preprocessing_distributions_comparison.png}{Comparaison des distributions avant et après prétraitement.}{fig:preprocessdist}

\projectfigure{08_preprocessing_boxplots_comparison.png}{Comparaison des boxplots selon les stratégies.}{fig:preprocessbox}

Les figures montrent que \texttt{log1p} réduit l'effet visuel des queues. StandardScaler seul change surtout le centre et l'échelle : il ne corrige pas intrinsèquement la skewness. De même, une transformation affine positive conserve la structure de l'IQR à un facteur d'échelle près. Le choix de B est donc motivé par la combinaison d'une compression logarithmique et d'une échelle commune, pas par une prétention à supprimer toute asymétrie.

\section{Implémentation réutilisable}
Le module \texttt{src/preprocessing.py} centralise la liste des features, construit les trois stratégies et sauvegarde les matrices ainsi que les objets de transformation. Cette organisation évite de recopier des formules dans plusieurs notebooks et permet de réappliquer les mêmes paramètres à de nouvelles observations.

Les fichiers exportés sont les suivants :
\begin{itemize}
  \item \path{customers_standard_scaled.csv} ;
  \item \path{customers_log_standard_scaled.csv} ;
  \item \path{customers_log_robust_scaled.csv}.
\end{itemize}
Les transformateurs sont sauvegardés dans \path{preprocessing_transformers.pkl}. Les trois matrices ont 440 lignes et six colonnes, sans valeur NaN ni infinie.

\section{Choix de la représentation principale}
La stratégie B, \texttt{log1p + StandardScaler}, est retenue comme candidate principale pour le clustering. Les stratégies A et C sont conservées comme alternatives de sensibilité. Ce choix est méthodologique et ne constitue pas un vainqueur universel : la qualité finale doit être jugée avec les partitions, les métriques et l'interprétation des profils.

\chapter{K-Means}

\section{Principe}
K-Means cherche une partition de $n$ observations en $K$ groupes. Chaque groupe est représenté par un centroïde, c'est-à-dire la moyenne des observations qui lui sont affectées. L'algorithme alterne une affectation au centroïde le plus proche et une mise à jour des centroïdes jusqu'à convergence.

L'initialisation \texttt{k-means++} choisit les centres de départ de manière plus informative que des tirages uniformes. Dans le notebook, \texttt{n\_init=20} lance plusieurs initialisations et conserve la solution de plus faible inertie. Le paramètre \texttt{random\_state=42} rend le résultat reproductible.

\begin{equation}
\mu_k=\frac{1}{|C_k|}\sum_{x_i\in C_k}x_i,
\qquad
J=\sum_{k=1}^{K}\sum_{x_i\in C_k}\lVert x_i-\mu_k\rVert^2.
\end{equation}

\section{Recherche de K}
Les valeurs $K=2$ à $K=10$ ont été testées sur les six features préparées. Le tableau \ref{tab:kmeanssearch} reprend l'inertie, la silhouette, le nombre d'itérations et les tailles observées.

\begin{table}[H]
\centering\scriptsize
\caption{Recherche de K pour K-Means.}
\label{tab:kmeanssearch}
\begin{tabular}{rrrrl}
\toprule
K & Inertie & Silhouette & Itérations & Effectifs \\
\midrule
2 & 1844,06 & 0,2903 & 13 & 252 ; 188 \\
3 & 1553,41 & 0,2583 & 7 & 145 ; 212 ; 83 \\
4 & 1386,86 & 0,1885 & 24 & 61 ; 129 ; 132 ; 118 \\
5 & 1266,15 & 0,1916 & 26 & 90 ; 96 ; 54 ; 100 ; 100 \\
6 & 1174,54 & 0,1953 & 19 & 91 ; 85 ; 110 ; 15 ; 56 ; 83 \\
7 & 1086,36 & 0,1958 & 20 & 20 ; 41 ; 48 ; 89 ; 84 ; 118 ; 40 \\
8 & 1024,05 & 0,1915 & 32 & 78 ; 28 ; 107 ; 46 ; 12 ; 21 ; 80 ; 68 \\
9 & 973,14 & 0,1965 & 12 & 19 ; 101 ; 84 ; 47 ; 64 ; 42 ; 28 ; 11 ; 44 \\
10 & 932,00 & 0,1905 & 13 & 18 ; 114 ; 16 ; 45 ; 40 ; 31 ; 42 ; 71 ; 11 ; 52 \\
\bottomrule
\end{tabular}
\end{table}

\projectfigure{09_kmeans_elbow.png}{Évolution de l'inertie selon le nombre de clusters.}{fig:elbow}
\projectfigure{10_kmeans_silhouette.png}{Évolution de la silhouette selon le nombre de clusters.}{fig:kmeanssil}

L'inertie décroît lorsque $K$ augmente, ce qui est attendu. Le coude n'est pas utilisé seul comme preuve d'une structure unique. La silhouette maximale parmi les valeurs testées est obtenue pour $K=2$. Cette solution présente également deux groupes de taille raisonnable, contrairement à certaines solutions plus fragmentées.

\section{Modèle final et profils}
Le modèle final est donc K-Means avec $K=2$. Les paramètres retenus sont :
\begin{itemize}
  \item \texttt{init=k-means++} ;
  \item \texttt{n\_init=20} ;
  \item \texttt{random\_state=42}.
\end{itemize}
Il obtient une inertie de 1844,0641 et une silhouette de 0,2903. Le cluster 0 rassemble 252 clients (57,27\%) et le cluster 1 rassemble 188 clients (42,73\%).

Sur les dépenses originales, le cluster 0 présente une moyenne Fresh de 13973,13 et des moyennes plus faibles pour Milk (2401,75), Grocery (2918,71) et Detergents\_Paper (491,94). Le cluster 1 est davantage orienté Milk (10346,36), Grocery (14697,06), Detergents\_Paper (6084,51) et Delicassen (2177,43), avec une moyenne Fresh de 9355,87.

\projectfigure{11_kmeans_cluster_profiles_heatmap.png}{Profils K-Means relatifs à la moyenne globale.}{fig:kmeansheat}
\projectfigure{12_kmeans_cluster_means.png}{Dépenses moyennes originales par cluster K-Means.}{fig:kmeansmeans}

Les numéros de clusters sont arbitraires. Ce qui est interprétable est la différence entre les profils, pas le numéro 0 ou 1 en lui-même.

\section{Projection PCA}
Une PCA à deux composantes a été appliquée uniquement pour visualiser les labels obtenus dans l'espace préparé. PC1 explique 44,08\% de la variance et PC2 27,19\%, soit 71,27\% cumulés.

\projectfigure{13_kmeans_pca_projection.png}{Projection PCA 2D colorée par les labels K-Means.}{fig:pcakmeans}

La figure est une projection et peut masquer des informations présentes dans les six dimensions. K-Means n'a pas été entraîné sur PC1 et PC2.

\section{Discussion}
Une silhouette de 0,2903 est positive mais modeste. Elle indique une organisation partielle des clients, avec des recouvrements possibles entre groupes. K-Means est retenu comme compromis pratique et descriptif, non comme preuve que le portefeuille possède deux catégories naturelles parfaitement séparées.

\chapter{Classification ascendante hiérarchique}

\section{Principe de la CAH}
La Classification Ascendante Hiérarchique commence avec une observation par groupe puis fusionne progressivement les groupes les plus proches. Le processus construit une hiérarchie représentée par un dendrogramme. Une coupe horizontale du dendrogramme détermine ensuite le nombre de groupes.

Le choix de la liaison modifie le résultat. Le saut minimum peut produire un effet de chaîne, le saut maximum favorise des groupes compacts, et la liaison moyenne utilise une distance moyenne. Ward choisit les fusions qui entraînent la plus faible augmentation de la variance intra-groupe. Dans ce projet, \texttt{linkage=ward} est appliqué aux six variables préparées.

\section{Recherche de la coupe}
Les valeurs K=2 à K=6 ont été comparées :
\begin{table}[H]
\centering
\caption{Scores CAH Ward pour les valeurs de K testées.}
\label{tab:cahsearch}
\begin{tabular}{rrrr}
\toprule
K & Silhouette & Taille minimale & Effectifs \\
\midrule
2 & 0,2585 & 178 & 178 ; 262 \\
3 & 0,2547 & 53 & 262 ; 53 ; 125 \\
4 & 0,2021 & 53 & 53 ; 208 ; 125 ; 54 \\
5 & 0,2025 & 18 & 208 ; 18 ; 125 ; 54 ; 35 \\
6 & 0,1707 & 18 & 125 ; 18 ; 150 ; 54 ; 35 ; 58 \\
\bottomrule
\end{tabular}
\end{table}

\projectfigure{14_cah_dendrogramme_ward.png}{Dendrogramme tronqué de la CAH avec liaison Ward.}{fig:dendro}
\projectfigure{15_cah_silhouette.png}{Silhouette selon K pour la CAH Ward.}{fig:cahsil}

Le meilleur score de silhouette est obtenu pour $K=2$, avec 0,2585. Les effectifs sont 178 et 262, donc la partition couvre les 440 clients sans bruit et sans groupe minuscule. Le score diminue pour les valeurs supérieures de K, même si l'inertie hiérarchique n'est pas utilisée comme critère principal dans ce notebook.

\section{Différence avec K-Means}
K-Means et la CAH Ward retiennent tous deux deux groupes, mais les affectations ne sont pas identiques. K-Means optimise directement l'inertie autour de centroïdes ; la CAH conserve la logique de fusions successives. Les profils moyens originaux permettent de comparer les tendances sans supposer que le cluster 0 d'une méthode correspond au cluster 0 de l'autre.

La CAH fournit un support visuel intéressant grâce au dendrogramme. Son score inférieur à K-Means dans cette représentation ne signifie pas que la CAH est inutile : elle décrit une hiérarchie et donne une autre lecture de la structure du portefeuille.

\chapter{DBSCAN}

\section{Principe basé sur la densité}
DBSCAN recherche des zones où les observations sont suffisamment denses. Le paramètre \texttt{eps} définit le rayon du voisinage et \texttt{min\_samples} le nombre minimal de points requis pour considérer une observation comme un point central (\textit{core point}). Les points proches d'un point central peuvent devenir des points de bord (\textit{border points}). Les observations qui ne rejoignent aucune région dense reçoivent le label -1 et sont dites bruit.

Cette logique est différente de K-Means : DBSCAN ne demande pas de fixer K à l'avance et peut reconnaître du bruit. En revanche, ses résultats sont sensibles à \texttt{eps}, \texttt{min\_samples}, à la dimension et à la représentation des données.

\section{Configurations testées}
Le notebook a testé \texttt{eps} entre 0,30 et 1,20, ainsi que \texttt{min\_samples} égal à 5 ou 10. Pour chaque configuration, le nombre de clusters, le nombre de points bruités, la couverture et la silhouette sur les points non bruités ont été calculés.

La configuration analysée dans la comparaison est \texttt{eps=0,80}, \texttt{min\_samples=5}. Elle produit trois clusters, dont les tailles sont 200, 5 et 5, ainsi que 230 points de bruit. La couverture est donc de 210/440, soit 47,73\%. La silhouette calculée sur les 210 points non bruités est 0,2885.

\begin{table}[H]
\centering
\caption{Résultats DBSCAN pour la configuration analysée et une configuration plus sélective.}
\label{tab:dbscan}
\scriptsize
\begin{tabularx}{\textwidth}{lrrrrX}
\toprule
Configuration & Clusters & Bruit & Couverture & Points évalués & Silhouette \\
\midrule
eps=0,80, min=5 & 3 & 52,27\% & 47,73\% & 210 & 0,2885 \\
eps=0,80, min=10 & 2 & 75,91\% & 24,09\% & 106 & 0,5372 \\
\bottomrule
\end{tabularx}
\end{table}

\projectfigure{16_dbscan_grid_results.png}{Taux de bruit et silhouette hors bruit selon les configurations DBSCAN.}{fig:dbscan}

\section{Pourquoi la silhouette élevée n'est pas suffisante}
La configuration \texttt{eps=0,80}, \texttt{min\_samples=10} obtient 0,5372 sur seulement 106 clients. La valeur est intéressante localement, mais 334 clients sont alors exclus, soit 75,91\% de bruit. Il serait trompeur de la présenter comme une meilleure segmentation du portefeuille complet.

La configuration analysée avec \texttt{min\_samples=5} conserve 210 points, mais elle produit deux groupes de cinq clients seulement et un groupe de 200 clients. Elle met en évidence des zones denses et des observations éloignées, mais ne fournit pas une partition globale équilibrée.

\begin{precaution}
Le bruit DBSCAN n'est pas automatiquement une anomalie commerciale. Il peut correspondre à un client atypique, à un mélange de profils ou à une densité insuffisante dans l'espace choisi. Une validation métier serait nécessaire avant toute décision.
\end{precaution}

\section{Conclusion DBSCAN}
Dans les configurations étudiées, aucune solution ne satisfait simultanément au moins deux clusters, une couverture d'au moins 70\% et une silhouette calculable satisfaisante. DBSCAN est donc utile comme analyse de densité et de cas atypiques, mais il n'est pas retenu comme partition principale du portefeuille.

\chapter{Comparaison des algorithmes}

\section{Cadre d'évaluation}
Une évaluation interne examine la structure produite sans disposer d'une vérité terrain. La silhouette combine cohésion et séparation. K-Means possède en plus l'inertie, directement liée à son objectif. Pour comparer correctement les méthodes, il faut aussi regarder les effectifs, la couverture et la facilité d'interprétation.

K-Means et CAH affectent chacun les 440 clients à un cluster. DBSCAN peut affecter des clients au bruit. Son score de silhouette est donc accompagné du nombre de points évalués et du taux de couverture.

\section{Tableau comparatif}
\begin{table}[H]
\centering\scriptsize
\caption{Comparaison finale des trois méthodes.}
\label{tab:comparison}
\begin{tabularx}{\textwidth}{lrrrrX}
\toprule
Méthode & Clusters & Silhouette & Points & Bruit & Effectifs et lecture \\
\midrule
K-Means & 2 & 0,2903 & 440 & 0\% & 252 ; 188. Partition complète, simple et équilibrée. \\
CAH Ward & 2 & 0,2585 & 440 & 0\% & 178 ; 262. Partition complète et dendrogramme lisible. \\
DBSCAN & 3 & 0,2885* & 210* & 52,27\% & 200 ; 5 ; 5, avec 230 points bruit. Partition globale déséquilibrée. \\
\bottomrule
\end{tabularx}
\end{table}

\textit{*} Le score DBSCAN est calculé sans les 230 points bruités. Sa couverture est 47,73\% ; il ne doit donc pas être comparé directement aux scores calculés sur les 440 observations.

\projectfigure{17_comparison_silhouette.png}{Comparaison visuelle des scores de silhouette.}{fig:comparisonsil}

\section{Analyse des avantages et limites}
K-Means obtient le meilleur compromis numérique parmi les solutions complètes testées. Il est facile à réexécuter, ses effectifs sont raisonnables et ses profils sont directement exploitables. Sa limite est la nécessité de choisir K et sa préférence pour une géométrie de groupes compacts.

La CAH donne une silhouette plus faible, mais elle permet d'observer la hiérarchie des fusions. Elle est pertinente lorsque la question porte sur plusieurs niveaux de regroupement. Son résultat dépend cependant de la liaison et de la coupe du dendrogramme.

DBSCAN est le seul à expliciter un bruit. Cette propriété est précieuse lorsqu'il existe des observations isolées ou des groupes de densité différente. Ici, le bruit représente toutefois plus de la moitié des clients dans la configuration analysée. Une silhouette hors bruit élevée dans une configuration très peu couverte ne suffit donc pas pour recommander DBSCAN comme segmentation globale.

\section{Choix final}
K-Means est retenu comme meilleur compromis \emph{parmi les configurations étudiées}. Ce choix ne signifie pas qu'il serait optimal pour toute autre représentation, tout autre dataset ou toute autre finalité. La silhouette de 0,2903 reste modeste et signale que les deux groupes se recouvrent partiellement.

\chapter{Interprétation commerciale prudente}

\section{Profils K-Means}
Le premier profil K-Means, de 252 clients, est caractérisé par une moyenne Fresh élevée (13973,13) et des niveaux plus faibles de Milk, Grocery et Detergents\_Paper. Le second profil, de 188 clients, présente des dépenses moyennes beaucoup plus élevées pour Milk (10346,36), Grocery (14697,06), Detergents\_Paper (6084,51) et Delicassen (2177,43). Il s'agit d'une description statistique des groupes, pas d'une étiquette commerciale officielle.

\section{Profils CAH Ward}
La CAH produit un groupe de 178 clients dont les moyennes Milk, Grocery et Detergents\_Paper sont élevées : respectivement 9338,17, 14247,11 et 6214,91. L'autre groupe, de 262 clients, est davantage caractérisé par Fresh (15016,10) et Frozen (4187,85), avec des niveaux plus faibles pour Milk, Grocery et Detergents\_Paper. Ces tendances ressemblent partiellement à la partition K-Means, mais les groupes ne sont pas identiques.

\section{Profils DBSCAN}
Dans la configuration \texttt{eps=0,80}, \texttt{min\_samples=5}, le groupe DBSCAN de 200 clients présente des moyennes intermédiaires, tandis que deux groupes de cinq clients ont des dépenses Milk, Grocery et Detergents\_Paper très élevées. Les 230 points de bruit ont un profil moyen qui mélange plusieurs comportements. Ce mélange rappelle que le label bruit indique une position dans la géométrie de densité, et non un type de client reconnu par le métier.

\projectfigure{18_profiles_methods_original.png}{Comparaison des profils moyens dans les unités originales.}{fig:profilesall}

\section{Rôle de Channel et Region}
Les variables Channel et Region n'ont pas été utilisées pour entraîner les distances. Elles peuvent néanmoins aider à décrire la composition des groupes après coup. Dans l'EDA, Channel 1 compte 298 clients et Channel 2 en compte 142. Les différences de moyennes sont particulièrement visibles pour Milk, Grocery et Detergents\_Paper. Cette information est cohérente avec l'existence de comportements d'achat différents selon le canal, mais elle ne doit pas être transformée en causalité.

Region est très déséquilibrée : la modalité 3 compte 316 clients contre 77 pour la modalité 1 et 47 pour la modalité 2. Une répartition observée dans un cluster doit donc être comparée à cette base globale et non interprétée isolément.

\section{Pistes commerciales, sous réserve de validation}
À titre d'hypothèses à valider avec les responsables métier, un groupe orienté Fresh et Frozen pourrait faire l'objet d'une analyse de disponibilité et de rotation des produits frais et surgelés. Un groupe orienté Milk, Grocery et Detergents\_Paper pourrait être étudié pour des offres liées aux achats récurrents d'épicerie et d'entretien. Ces propositions ne sont pas des résultats démontrés par le clustering : elles doivent être confrontées aux marges, aux fréquences de commande, aux coûts logistiques et à l'historique temporel.

\begin{precaution}
Un cluster décrit une proximité dans les six dépenses choisies. Il ne permet pas de conclure sur la rentabilité, la fidélité, la cause d'achat ou l'efficacité d'une campagne marketing. Ces dimensions nécessiteraient des variables supplémentaires et une validation externe.
\end{precaution}

\chapter{Conclusion générale}

Le projet a étudié la segmentation de 440 clients d'un distributeur en gros à partir de six dépenses annuelles. La démarche a commencé par la conservation du CSV original, le contrôle de sa qualité et une EDA complète. Les données ne contiennent aucune valeur manquante, aucun doublon exact et aucune dépense négative. Les distributions sont positivement asymétriques et la règle IQR signale 172 valeurs cumulées, sans justifier leur suppression automatique.

Trois stratégies de prétraitement ont été comparées. La stratégie principale \texttt{log1p + StandardScaler} compresse les grandes valeurs puis donne aux six dépenses une échelle comparable. StandardScaler seul et \texttt{log1p + RobustScaler} restent disponibles pour l'analyse de sensibilité.

K-Means, avec \texttt{k-means++}, \texttt{n\_init=20} et \texttt{random\_state=42}, retient $K=2$. Sa silhouette est 0,2903, son inertie est 1844,0641 et ses groupes contiennent 252 et 188 clients. La CAH Ward retient également $K=2$, avec une silhouette de 0,2585 et des effectifs 178 et 262. DBSCAN met en évidence une structure de densité, mais la configuration analysée produit 52,27\% de bruit ; sa silhouette 0,2885 est calculée sur 210 clients seulement.

K-Means est donc retenu comme meilleur compromis parmi les configurations testées, car il fournit la meilleure silhouette sur l'ensemble complet avec des effectifs équilibrés. Cette conclusion reste prudente : 0,2903 est une séparation positive mais modeste. Les groupes ne doivent pas être présentés comme parfaitement séparés ni comme des catégories causales.

Les perspectives sont l'ajout de variables métier, l'étude temporelle des commandes, la validation par des experts commerciaux et l'analyse de stabilité sur d'autres périodes ou échantillons. Une future comparaison pourrait également étudier systématiquement la sensibilité des résultats aux trois prétraitements, sans remettre en cause la conservation du dataset original.


\appendix
\chapter{Compléments et repères pour la soutenance}

\section{Paramètres principaux}
\begin{table}[H]
\centering
\caption{Paramètres utilisés dans les notebooks de clustering.}
\label{tab:params}
\begin{tabularx}{\textwidth}{l l X}
\toprule
Méthode & Paramètres & Interprétation \\
\midrule
K-Means & $K=2$, \texttt{init=k-means++}, \texttt{n\_init=20}, \texttt{random\_state=42} & Deux groupes autour de centroïdes. \\
CAH & $K=2$, \texttt{linkage=ward} & Coupe de la hiérarchie donnant la meilleure silhouette testée. \\
DBSCAN & \texttt{eps=0,80}, \texttt{min\_samples=5} & Trois groupes et un label bruit pour 230 clients. \\
\bottomrule
\end{tabularx}
\end{table}

\section{Formules utiles}
\begin{align}
\text{log1p}(x)&=\log(1+x),\\
z&=\frac{x-\mu}{\sigma},\\
IQR&=Q_3-Q_1,\\
s(i)&=\frac{b(i)-a(i)}{\max(a(i),b(i))}.
\end{align}

\section{Lecture rapide des fichiers Python}
\texttt{src/check\_dataset.py} charge le CSV et vérifie dimensions, colonnes, types, valeurs manquantes, doublons et dépenses attendues. \texttt{src/preprocessing.py} centralise l'ordre des six dépenses et construit les pipelines de transformation. Les notebooks utilisent ensuite ces fichiers préparés sans écraser la source.

\section{Glossaire}
\begin{description}
\item[Centroïde] moyenne vectorielle des observations d'un cluster K-Means.
\item[Couverture] proportion des observations effectivement affectées à un cluster ; pour DBSCAN, elle exclut le bruit.
\item[Inertie] somme des distances quadratiques aux centroïdes, critère spécifique à K-Means.
\item[Silhouette] mesure interne combinant proximité au propre cluster et séparation du cluster voisin.
\item[Bruit] observation DBSCAN qui n'est rattachée à aucune zone dense selon les paramètres choisis.
\item[Skewness] mesure d'asymétrie d'une distribution.
\end{description}

\section{Ressources graphiques}
Les 18 figures utilisées dans le rapport sont copiées dans \texttt{reports/latex/figures/}. Les dix tableaux CSV exportés sont conservés dans \texttt{reports/latex/tables/} pour consultation et réutilisation dans Overleaf ou dans une analyse ultérieure.


\backmatter
\printbibliography[heading=bibintoc,title={Bibliographie}]
\end{document}
